\chapter{Introducción} \label{ch:introduccion}

% ============================================================
%  BLOQUE 1 — CONTEXTO: RELEVANCIA DEL FENÓMENO Y ESTADO DEL ARTE
% ============================================================

% --- 1.1: Rol climático y dinámico global; algoritmos de seguimiento ---
Los ciclones extratropicales (ETC, por sus siglas en inglés) son sistemas sinópticos recurrentes en latitudes medias y altas que dominan el tiempo severo de estas regiones.
\citet{hoskins2005new} caracterizan los \textit{storm tracks} del HS, los corredores preferenciales por donde estos ciclones canalizan gran parte del transporte latitudinal de energía del hemisferio.
La posibilidad de cuantificar sistemáticamente estos sistemas se consolidó con el desarrollo de algoritmos de seguimiento objetivo. \citet{sinclair1994objective} propuso el uso de la vorticidad relativa para identificar y rastrear ciclones independientemente del flujo de fondo, mientras que \citet{simmonds1999southern} desarrollaron un esquema complementario, basado en la Laplaciana de la presión a nivel del mar, para el seguimiento automático de sistemas en el HS; más recientemente, \citet{padilhareinke2024objective} refinaron específicamente el enfoque de vorticidad relativa para la región.

% --- 1.2: Estado del arte regional — climatología, diversidad estructural y regionalización ---
Inicialmente, autores como \citet{gan1991surface} cuantificaron que la ciclogénesis en la región sudamericana es más frecuente en invierno y que su variabilidad interanual repercute en la precipitación del sur de Brasil. 
Poco después, la climatología hemisférica de \citet{jones1993climatology} confirmó a mayor escala esta señal, identificando a la región sudamericana como un núcleo de densidad ciclónica propio de las latitudes extratropicales, distinto de las concentraciones mayores en latitudes cercanas al polo.
Sobre el Atlántico Sur occidental, esta actividad ciclónica es particularmente intensa, con aproximadamente 200 ciclones extratropicales registrados por año \citep{reboita2026meteorology}.
Ya en un trabajo más reciente, \citet{reboita2010south} construyeron una climatología para el Atlántico Sur e identificaron tres centros preferenciales de ciclogénesis: uno al sur, vinculado a la inestabilidad baroclínica de los oestes, y dos más al norte, próximos a Uruguay y el sur de Brasil, vinculados a la ciclogénesis de sotavento inducida por la onda topográfica de los Andes.
Trabajos posteriores, como \citet{gramcianinov2020analysis} y recientemente \citet{padilhareinke2026characterization}, han precisado la distribución espacio-temporal de estos sistemas, aportando conocimiento sobre la frecuencia y localización de la actividad ciclónica en la región.
Sin embargo, esta actividad ciclónica no es dinámicamente uniforme. \citet{hart2003cyclone} demostró que los sistemas transitan a lo largo de un continuo estructural y no de categorías estáticas, y en el Atlántico Sur esta heterogeneidad se manifiesta en ciclones de estructura térmica híbrida \citep{evans2012climatology} y en sistemas que siguen el modelo de Shapiro-Keyser, con fractura del frente frío y seclusión cálida en la madurez \citep{gozzo2013air}.
El Atlántico Sur se aparta con frecuencia de la conceptualización clásica de ciclón, tal como lo evidencian las variaciones morfológicas respaldadas por una climatología regional de largo plazo \citep{gozzo2014subtropical}.
Esta diversidad estructural coexiste, además, con una regionalización dinámica clara. La Cordillera de los Andes, cuyo forzante orográfico incrementa la baroclinicidad a sotavento \citep{gan1994influence}, combinada con los gradientes de temperatura superficial del mar \citep{reboita2026meteorology}, consolida focos de actividad ciclogenética diferenciados en la costa este de Sudamérica, cuyos mecanismos físicos específicos se desarrollan en la Sección~\ref{sec:tb_regiones} de los Fundamentos.
\label{sec:definicion_regiones}%
Dada esta heterogeneidad estructural y regional, este estudio se centra en tres regiones de ciclogénesis documentadas por diversos autores \citep{crespo2020potential, gramcianinov2019properties, reboita2026meteorology} y delimitadas siguiendo la propuesta de \citet{gramcianinov2019properties}: la Sur-Sudeste de Brasil (\textit{Southeast Brazil}, SBR), la cuenca de descarga del Río de la Plata (\textit{La Plata Basin}, LPB) y la costa sureste de Argentina (\textit{Argentina}, ARG).

% --- 1.3: Mecanismos dinámicos internos ya establecidos ---
Al este de los Andes, la convergencia de viento en niveles bajos favorece las condiciones que producen precipitaciones intensas en el sureste de Sudamérica, cuya liberación de calor latente retroalimenta y refuerza la intensificación del sistema \citep{vera2002cold}.
El transporte meridional de humedad desde los trópicos hacia estos sistemas es mediado por el Chorro de Capas Bajas de Sudamérica \citep{reboita2026meteorology}, mientras que los flujos turbulentos de calor superficiales océano-atmósfera, como demuestran \citet{reboita2022from}, no solo inestabilizan la atmósfera, sino que actúan como fuente diabática que refuerza la intensidad del sistema.
Esta advección cálida se organiza en torno a estructuras de mesoescala conocidas como cintas transportadoras, la Cinta Transportadora Cálida y la Cinta Transportadora Fría (WCB y CCB, respectivamente), cuyo papel en la distribución espacial de viento y precipitación se desarrolla en la Sección~\ref{subsec:tb_modelos}.
En esta misma línea, \citet{gentile2025response} anticipan, mediante modelaciones de alta resolución, que el sector cálido de los ETC más intensos se intensificará en un clima futuro, incrementando significativamente los vientos y la precipitación durante sus etapas de desarrollo.
El viento en superficie y la precipitación son, por ello, las variables elegidas para caracterizar estos impactos, según se detalla en la Sección~\ref{subsec:tb_compound_def}.

% ============================================================
%  BLOQUE 2 — PROBLEMA: EL VACÍO EN LA CARACTERIZACIÓN ESTRUCTURAL
% ============================================================

Pese a este conocimiento consolidado sobre dónde, cuándo y por qué mecanismo se forman los ETC del Atlántico Sur, persiste un problema no resuelto en \emph{cómo} se distribuyen espacialmente sus impactos de superficie a lo largo del ciclo de vida.
\citet{corner2025classification} muestran que la intensidad de un ciclón no puede capturarse con una única métrica, dado que las relaciones entre distintas medidas de intensidad, incluida la de los vientos en superficie, no son lineales entre sí, y \citet{sinclair2023relationship} confirman este límite al advertir que un ETC puede ser dinámicamente intenso pero generar poca lluvia si carece del suministro de humedad adecuado.
A esta insuficiencia escalar se suma una insuficiencia posicional. Los máximos de viento no se distribuyen uniformemente alrededor del centro del ciclón, sino que se concentran en cuadrantes específicos cuya ubicación varía con estructuras de mesoescala y con la fase del ciclo de vida \citep{cardoso2022synoptic, eisenstein2023identification}, mecanismo que se desarrolla en la Sección~\ref{subsec:tb_kde_asimetria} de los Fundamentos.
Esta insuficiencia se confirma también en el HS, donde \citet{pepler2020dimensional} muestran, para el sureste de Australia, que los ciclones sustentados en niveles verticales altos presentan vientos e intensidad de lluvia sistemáticamente mayores que los ciclones configurados solo en niveles bajos.
En Sudamérica, los estudios de caso de \citet{gozzo2013air} y \citet{reboita2022from} han revelado, en este mismo sentido, que los ciclones regionales pueden exhibir evoluciones que se apartan de los modelos tradicionales, lo que confirma que la distribución espacial de sus campos de superficie durante la evolución del sistema no puede darse por establecida a priori.
Aunque el viento y la precipitación extremos responden a mecanismos físicos y escalas temporales distintos, \citet{chen2025characteristics} muestran que su coocurrencia, si bien relativamente poco frecuente, no es despreciable, por lo que caracterizar la morfología superficial del sistema exige analizar ambas variables de forma conjunta y con una descripción de su covariabilidad estructural a lo largo del ciclo de vida, no solo de su organización espacial individual.
\citet{han2025system} enfatizan, en consecuencia, que se requieren marcos de clasificación que vinculen explícitamente la evolución del sistema con sus manifestaciones de viento y precipitación, necesidad que campañas de observación aerotransportadas han vuelto más evidente al confirmar la presencia de convección interna invisible para los productos de resolución convencional \citep{kaylee2025convection}.
Esta caracterización estructural ha estado, además, limitada por la resolución espacial de los datos disponibles. \citet{priestley2022improved} demuestran que los modelos de baja resolución subestiman los vientos extremos asociados a los ciclones al no resolver explícitamente sus estructuras de mesoescala, y \citet{neu2013intercomparison} señalan que la variabilidad en forma y tamaño de los ETC genera divergencias sistemáticas entre métodos de detección, particularmente pronunciadas en regiones de topografía compleja.

% ============================================================
%  BLOQUE 3 — JUSTIFICACIÓN E IMPLICACIONES
% ============================================================

Resolver este vacío tiene una relevancia que trasciende lo descriptivo.
En el HS, los ETC modulan el oleaje regional y representan una amenaza costera recurrente \citep{gramcianinov2023impact}, y \citet{sasaki2021intraseasonal} muestran que esta modulación presenta variabilidad significativa a escala intraestacional en el Atlántico Sur occidental, con un período dominante de aproximadamente 40 días asociado a cambios en el viento en superficie, lo que sugiere que el riesgo costero puede modularse en escalas temporales que exceden la duración de un evento sinóptico individual.
Esta amenaza no es estática. Los registros históricos de 40 años ya muestran una tendencia significativa hacia ciclones de mayor tamaño e intensidad en el HS \citep{simmonds2000variability}, tendencia que las proyecciones de calentamiento futuro sugieren se intensificará específicamente en la componente de precipitación, con una tasa de escalamiento cercana al 7\%/K \citep{kodama2019perspective}.
En el litoral sur de Brasil, \citet{bartolomei2024extremos} documentan la letalidad y los impactos socioeconómicos causados por ciclones extratropicales invernales, mientras que \citet{miranda2026patterns} confirman que estos sistemas son los principales motores de tormentas costeras en el litoral sur de Brasil; la asociación directa entre el viento observado en la costa y estos sistemas confirma, además, la ocurrencia de pérdidas socioeconómicas asociadas \citep{bitencourt2010relating}.
Una caracterización de dónde y en qué fase del ciclo de vida se concentran los máximos de viento y precipitación, en lugar de depender de una métrica escalar de intensidad central, es, por lo tanto, la información que permitiría anticipar con mayor precisión la localización y magnitud de estos impactos costeros, y contribuiría con evidencia regional al esfuerzo internacional, señalado por \citet{han2025system}, de construir marcos de clasificación objetiva que vinculen explícitamente la evolución del sistema con sus manifestaciones de superficie.
La viabilidad de emprender esta caracterización reside en tres desarrollos recientes que, combinados, la hacen posible por primera vez para el Atlántico Sur.
Primero, la disponibilidad del reanálisis ERA5 \citep{hersbach2020era5}, cuya resolución espacial ($\sim$31 km) y temporal (horaria) permite resolver los gradientes de presión y los flujos de humedad que generaciones anteriores de datos suavizaban; validaciones específicas para la región \citep{gramcianinov2020analysis} y comparaciones frente a otras fuentes de datos para la detección de ciclogénesis en Sudamérica \citep{dalanhese2023new} respaldan esta elección, aunque ERA5 presenta sesgos conocidos en la representación de los extremos más intensos, discutidos en la Sección~\ref{sec:datos_era5}.
Segundo, la base de datos de trayectorias de \citet{coutodesouza2024new}, fundamentada en el seguimiento por vorticidad relativa a 850 hPa ($\zeta_{850}$) en lugar de la presión al nivel del mar, cuyos fundamentos dinámicos y validación para el HS se desarrollan en la Sección~\ref{sec:tb_vorticidad}, y que además incorpora la segmentación objetiva del algoritmo \textit{CycloPhaser}, la cual define cuatro fases dinámicas discretas, incipiente, intensificación, madurez y decaimiento, y supera las limitaciones de las clasificaciones clásicas en el contexto del Atlántico Sur \citep{deSouza2025CycloPhaser}, en sintonía con el marco de referencia que ofrece el modelo Shapiro-Keyser \citep{shapiro1990life,shapiro1999bridge} para interpretar la distribución espacial de los máximos de viento y precipitación (Sección~\ref{subsec:tb_modelos}).
Tercero, la disponibilidad de técnicas estadísticas capaces de capturar simultáneamente la asimetría estructural de cada campo y su covariabilidad. Dado que la estructura del ciclón presenta una variabilidad espacial continua en la que la circulación del viento y los núcleos de precipitación pueden covariar durante la evolución del sistema, el análisis de Funciones Ortogonales Empíricas Multivariadas (MEOF) permite extraer modos de variabilidad inherentemente acoplados entre variables físicamente distintas \citep{liang2018multivariate}, a diferencia del análisis univariado de EOF aplicado por separado a cada campo, y captura así la organización espacial conjunta de ambos.
La utilidad de estos patrones y de sus componentes principales trasciende la caracterización descriptiva, con aplicación práctica incluso en el pronóstico de precipitación \citep{sun2022evaluation}, y su robustez estadística puede evaluarse mediante técnicas de remuestreo (\textit{bootstrap}) que distinguen las señales físicas consistentes del ruido de muestreo propio de cada subconjunto.

% ============================================================
%  BLOQUE 4 — PREGUNTA DE INVESTIGACIÓN, HIPÓTESIS Y OBJETIVO
% ============================================================

La combinación de estos tres desarrollos, los datos de ERA5, la detección dinámica por $\zeta_{850}$ y la clasificación objetiva de fases mediante \textit{CycloPhaser}, ya disponibles en la base de \citet{coutodesouza2024new}, junto con el análisis multivariado (MEOF) desarrollado en esta investigación, constituye la base técnica para cuantificar la evolución de la estructura espacial de los campos de viento y precipitación en la región, y para construir, a partir de las fases previamente clasificadas, prototipos estructurales que sinteticen la morfología típica del ciclón para cada región de génesis y fase evolutiva.
Con este marco de datos, detección dinámica y herramientas analíticas, se plantea la pregunta que guía esta investigación: ¿cómo evoluciona espacialmente la distribución de viento y precipitación en los ETC del Atlántico Sur, y qué configuraciones caracterizan esta organización durante las fases incipiente, intensificación, madurez y decaimiento?
Se postula que los campos de viento y precipitación presentan patrones de asimetría sistemática que varían entre fases evolutivas, entre regiones de génesis y entre estratos de intensidad dinámica, contrastando la Población Global de ciclones frente al subconjunto de mayor vorticidad, y que dichos patrones no son capturables mediante métricas escalares de intensidad central.

\section{Objetivos del Proyecto de Investigación}

\subsection{Objetivo General}
Caracterizar la evolución de la estructura espacial del viento a 10 metros y precipitación durante las fases del ciclo de vida de ETC en el Atlántico Sur, caracterizando sus patrones de variabilidad y estimación de densidad probabilística de máximos, con estratificación por región de génesis.

\subsection{Objetivos Específicos}
\begin{enumerate}
    \item Consolidar una base de datos lagrangiana de los ETC del Atlántico Sur que asocie los campos horarios de alta resolución de viento y precipitación (ERA5) con las trayectorias y fases dinámicas (incipiente, intensificación, madurez y decaimiento), estratificando la muestra resultante por región de génesis (SBR, LPB y ARG) y por intensidad dinámica.

    \item Cuantificar la distribución estadística y la asimetría espacial de los máximos de viento a 10 m y precipitación mediante estimación de densidad por núcleos, caracterizando tanto la distribución de probabilidad de sus magnitudes (PDF unidimensional) como la localización preferencial de dichos máximos respecto al centro del ciclón (PDF espacial).

    \item Identificar los modos dominantes de variabilidad espacial de los campos de viento y precipitación mediante Funciones Ortogonales Empíricas univariadas (EOF) y multivariadas (MEOF), evaluando el patrón estructural individual de cada variable y la covariabilidad acoplada.

    \item Construir prototipos estructurales sintéticos por fase evolutiva y región de génesis, integrando los campos medios condicionados al modo dominante de covariabilidad (MEOF) con las densidades probabilísticas de localización de máximos (PDFe), para cuantificar la magnitud característica de viento y precipitación.
\end{enumerate}
